% from ega3-1-fr.tex
\oldpage[0\textsubscript{III}]{36}
de $\mathcal C$ dans $\mathcal C'$, permutant aux
limites inductives (on suppose que les limites inductives filtrantes existent
et sont exactes dans $\mathcal C$ et $\mathcal C'$). Alors
$\mathrm R^\bullet T(K^\bullet)$ est isomorphe à la limite inductive
\[
  \varinjlim_{r\to-\infty}\mathrm R^\bullet T(K_{(r)}^\bullet)
\]
lorsque $r$ tend vers $-\infty$.
\end{lemma}

La construction d'une résolution injective de Cartan-Eilenberg de $K^\bullet$
se fait en choisissant arbitrairement, pour chaque $i$, une résolution
injective $(X_{\mathrm B}^{ij})_{j\geq0}$ de $\mathrm B^i(K^\bullet)$ et une
résolution injective $(X_{\mathrm H}^{ij})_{j\geq0}$ de
$\mathrm H^i(K^\bullet)$~; cela fait, la méthode de construction montre que
la résolution injective $(L^{ij})_{j\geq0}$ de $K^i$ et les opérateurs de
dérivation $L^{i,j}\to L^{i+1,j}$ \emph{ne dépendent que des résolutions}
$(X_{\mathrm B}^{i,\bullet})$, $(X_{\mathrm H}^{i,\bullet})$ et
$(X_{\mathrm B}^{i+1,\bullet})$ (M, XVII, 1.2). Or, il est clair que l'on a
$\mathrm B^i(K_{(r)}^\bullet)=\mathrm B^i(K^\bullet)$ et
$\mathrm H^i(K_{(r)}^\bullet)=\mathrm H^i(K^\bullet)$ pour $i\geq r+1$. On
a, d'autre part, pour chaque $r$, une injection canonique
$\varphi_{r-1,r}^\bullet:K_{(r)}^\bullet\to K_{(r-1)}^\bullet$,
$\varphi_{r-1,r}^i$ étant l'identité pour $i\neq r-1$. La remarque précédente
montre que, si $L^{\bullet\bullet}=(L^{ij})$ est une résolution injective de
Cartan-Eilenberg de $K^\bullet$, on peut pour chaque $r$ définir une résolution
injective de Cartan-Eilenberg
$L_{(r)}^{\bullet\bullet}=(L_{(r)}^{ij})$ de $K_{(r)}^\bullet$ telle que
$L_{(r)}^{ij}=0$ pour $i<r$ et $L_{(r)}^{ij}=L^{ij}$ pour $i\geq r+1$. On peut
d'autre part définir un morphisme de bicomplexes
$\Phi_{r-1,r}^{\bullet\bullet}:L_{(r)}^{\bullet\bullet}\to
L_{(r-1)}^{\bullet\bullet}$ correspondant à $\varphi_{r-1,r}^\bullet$, et la
méthode de définition de ce morphisme (\emph{loc. cit.}) montre encore qu'on
peut le construire de sorte que $\Phi_{r-1,r}^{ij}$ soit l'identité pour
$i\neq r$ et $i\neq r-1$. On a ainsi défini un système inductif
$(L_{(r)}^{\bullet\bullet})$ de bicomplexes de $\mathcal C$, dont
$L^{\bullet\bullet}$ est évidemment la limite inductive lorsque $r$ tend vers
$-\infty$~; en raison de la permutabilité des sommes directes et des limites
inductives, le complexe simple associé à $L^{\bullet\bullet}$ est aussi limite
inductive du complexe simple associé à $L_{(r)}^{\bullet\bullet}$. Comme $T$
permute par hypothèse aux limites inductives et qu'il en est de même de la
cohomologie (en raison de l'exactitude du foncteur $\varinjlim$), on a bien
\[
  \mathrm H^\bullet(T(L^{\bullet\bullet}))=
  \varinjlim\mathrm H^\bullet(T(L_{(r)}^{\bullet\bullet}))
\]
à un isomorphisme près.

Le lemme (\hyperref[0.11.5.1-fr]{11.5.1}) permet d'étendre à des complexes
$K^\bullet$ quelconques, par passage à la limite inductive, des propriétés
valables pour les complexes \emph{limités inférieurement}. Comme premier
exemple, nous prouverons~:

\begin{proposition}[11.5.2]
\label{0.11.5.2-fr}
Sous les hypothèses de (\hyperref[0.11.5.1-fr]{11.5.1}) concernant
$\mathcal C$, $\mathcal C'$ et $T$, $\mathrm R^\bullet T(K^\bullet)$ est un
foncteur cohomologique dans la catégorie abélienne des complexes de
$\mathcal C$.
\end{proposition}

\begin{proof}
Montrons qu'on peut se ramener au cas des complexes limités inférieurement~:
si on a une suite exacte
\[
  0\longrightarrow K'^\bullet\longrightarrow K^\bullet
  \longrightarrow K''^\bullet\longrightarrow0
\]
de complexes, on en déduit évidemment pour chaque $r$ une suite exacte
\[
  0\longrightarrow K_{(r)}'^\bullet\longrightarrow K_{(r)}^\bullet
  \longrightarrow K_{(r)}''^\bullet\longrightarrow0,
\]
d'où par hypothèse, une suite exacte
\[
  \cdots\longrightarrow\mathrm R^nT(K_{(r)}'^\bullet)
  \longrightarrow\mathrm R^nT(K_{(r)}^\bullet)
  \longrightarrow\mathrm R^nT(K_{(r)}''^\bullet)
  \xrightarrow{\partial}\mathrm R^{n+1}T(K_{(r)}'^\bullet)
  \longrightarrow\cdots
\]
ces suites exactes formant un système inductif~; le lemme
(\hyperref[0.11.5.1-fr]{11.5.1}) et l'exactitude du foncteur $\varinjlim$
démontrent que l'on a une suite exacte
\[
  \cdots\longrightarrow\mathrm R^nT(K'^\bullet)
  \longrightarrow\mathrm R^nT(K^\bullet)
  \longrightarrow\mathrm R^nT(K''^\bullet)
  \xrightarrow{\partial}\mathrm R^{n+1}T(K'^\bullet)
  \longrightarrow\cdots
\]
\end{proof}

Pour traiter le cas des complexes limités inférieurement, on peut se borner à
ceux tels que $K^i=0$ pour $i<0$~; ils forment évidemment une catégorie
abélienne $\mathcal K$.

\begin{lemma}[11.5.2.1]
\label{0.11.5.2.1-fr}
Dans $\mathcal K$, soit $\mathfrak J$ l'ensemble des complexes
$Q^\bullet=(Q^i)_{i\geq0}$ ayant les
